\booksection{ملحق: تحليل توسيع ذاكرة هواتف Android}

\chapterintro{عند شراء هاتف، غالباً ما ترى إعلاناً مثل “ذاكرة 12GB+8GB” — وكأن 8GB من الذاكرة العشوائية أضيفت من العدم. لكن هل يمكنها حقاً جعل الهاتف أسرع؟ وهل قد تؤذي الهاتف؟ في هذا القسم سنوضح تقنية “توسيع الذاكرة” بالكامل.}

\booksubsection{أولاً: التعريف الجوهري}

“توسيع الذاكرة” في هواتف Android (يسمى أيضاً توسيع الذاكرة، دمج الذاكرة، RAM Plus، إلخ)\textbf{ليس إضافة فعلية لرقاقة الذاكرة العشوائية الفيزيائية}، بل تقنية برمجية على مستوى النظام: من خلال تحويل جزء من مساحة التخزين الداخلي للهاتف (ROM/UFS Flash) ديناميكياً، محاكاته كذاكرة عشوائية افتراضية (RAM) للاستخدام.

على سبيل التشبيه: RAM الفيزيائية هي “طاولة العمل السريعة”، والذاكرة الموسعة هي “غرفة التخزين المؤقتة”. عندما لا تكفي طاولة العمل، يقوم النظام بنقل البيانات غير المستخدمة مؤقتاً إلى غرفة التخزين، وعند الحاجة يعيدها. غرفة التخزين كبيرة، لكن النقل ذهاباً وإياباً أبطأ بالتأكيد من وضعها مباشرة على طاولة العمل.

\booksubsection{ثانياً: المبدأ التقني}

طريقة تنفيذ الشركات المصنعة الرئيسية نوعان:

\subsubsection*{ضغط zRAM}

تقسيم مساحة بيانات مضغوطة داخل RAM، تخزين رمز التطبيقات غير النشط مضغوطاً، نسبة الضغط حوالي 2:1. البيانات لا تزال في RAM، سريعة، لا تتلف الـ Flash، لكنها تحتاج إلى CPU للمشاركة في الضغط وفك الضغط.

\subsubsection*{قسم swap}

تقسيم مساحة ثابتة مباشرة في UFS Flash كذاكرة افتراضية، كتابة بيانات العمليات الخلفية إلى التخزين. يعتمد على عرض نطاق التخزين، سرعة القراءة والكتابة أقل بكثير من الذاكرة الفيزيائية.

“ذاكرة 12GB+8GB” التي تعلن عنها الشركات المصنعة عادةً هي تركيبة 12GB RAM فيزيائية + 8GB ذاكرة افتراضية من الـ Flash. تجدر الإشارة إلى أن zRAM لا يحتل تخزين الجسم افتراضياً، فقط عندما يمتلئ zRAM بالكامل بالبيانات، يقوم النظام بإنشاء ملف تبادل على التخزين، وهذا نادر الحدوث.

\booksubsection{ثالثاً: فجوة السرعة هي الاختناق الرئيسي}

القيد الأساسي لهذه التقنية يكمن في اختلاف السرعة. UFS 4.0 Flash السائد حالياً سرعة قراءة تسلسلية حوالي 4200MB/s، أي 1/10–1/15 من ذاكرة LPDDR5X الفيزيائية. لذا في مرحلة ضغط zRAM منخفضة الحمل تكون السرعة قريبة من RAM الحقيقية، لكن بمجرد الدخول في سيناريوهات الحمل العالي، تنخفض السرعة إلى جزء من عشرات.

\booksubsection{رابعاً: سيناريوهات التطبيق وحدود الفائدة}

توسيع الذاكرة ليسحلاً سحرياً، التأثير يختلف حسب الطراز وسيناريو الاستخدام:

\begin{itemize}
  \item \textbf{الحالات المفيدة}: الطرازات بذاكرة فيزيائية 6GB + UFS 3.1 فأكثر، ومساحة متبقية $\geq$30GB، بعد تفعيل توسيع 3GB يمكن زيادة عدد التطبيقات النشطة في الخلفية بـ 3-4، انخفاض معدل إعادة التحميل بنسبة 68\%، مناسب لمستخدمي المهام المتعددة الثقيلين مثل مندوبي توصيل الطعام وسائقي سيارات الأجرة.
  \item \textbf{السيناريوهات اليومية}: تحسن ملحوظ في تشغيل WeChat متعدد، إبقاء تبويبات المتصفح نشطة، استجابة رسائل الاتصال الفوري، وقت التشغيل البارد للتطبيقات الخلفية يقل بوضوح.
  \item \textbf{السيناريوهات المحدودة}: تحسن محدود لتحميل الألعاب الكبيرة، تحرير فيديو 4K وغيرها من سيناريوهات النطاق الترددي العالي.
  \item \textbf{الحالات الضارة}: الطرازات بذاكرة 8GB + UFS 2.2 ومساحة متبقية غير كافية، بعد التفعيل تتسع تقلبات معدل الإطارات في الألعاب، استهلاك طاقة أكبر عند تصفح الفيديوهات القصيرة، تقطيع عند التبديل.
  \item \textbf{الحالات غير الفعالة}: الطرازات بذاكرة كبيرة 16GB فأكثر الفائدة شبه معدومة، على الأكثر إبقاء تطبيق واحد إضافي نشطاً، لكنها تزيد حمل CPU بنسبة 12\%، وترفع الحرارة 2-3℃. للطرازات بـ RAM أكبر من 12GB، يُنصح بإيقافها مباشرة لتجنيب حسابات الضغط غير ذات المعنى.
\end{itemize}

\booksubsection{خامساً: تسمية كل شركة وطريقة التفعيل}

تختلف التسميات بين العلامات التجارية: Xiaomi تسميه “توسيع الذاكرة”، Huawei تسميه “دمج الذاكرة”، Samsung تسميه “RAM Plus”، vivo تسميه “Extended RAM”، OnePlus تسميه “Memory Extension”. هذه الميزة قدرة أصلية للنظام، لا علاقة مباشرة لها بإصدار Android الأساسي.

مسار التفعيل العام: الدخول إلى [الإعدادات] [المزيد من الإعدادات] إيجاد خيار [الذاكرة] أو [الأداء] تفعيل “توسيع الذاكرة” واختيار مستوى السعة (عادةً 2GB–7GB)، يسري بعد إعادة التشغيل. بعد التفعيل سيعرض النظام تركيبة ذاكرة مثل “8GB+4GB”.

\begin{tcbwarn}{تحذيرات}
\begin{itemize}
  \item توسيع الذاكرة سيحتل جزءاً من مساحة تخزين الجسم، يرجى التأكد من أن التخزين المتبقي لا يقل عن 20-30GB، وإلا قد يؤثر على سلاسة النظام.
  \item استخدام Flash كذاكرة افتراضية سيكون له خسارة قراءة وكتابة، رغم أن عمر UFS Flash الحديث طويل، لكن swap عالي التردد على المدى الطويل قد يسرع شيخوخة التخزين.
  \item إذا كانت الذاكرة الفيزيائية للهاتف صغيرة (مثل 4GB وأقل)، فإن تفعيل توسيع الذاكرة يخفف فقط من المشكلة المؤقتة، ولا يحل مشكلة التقطيع من جذورها.
\end{itemize}

\end{tcbwarn}

\begin{tcbinfo}{خلاصة}
توسيع الذاكرة تقنية تحسين ناضجة على مستوى النظام، يمكنها تحسين قدرات المهام المتعددة ومعدل إبقاء التطبيقات نشطة في الخلفية للطرازات المتوسطة والمنخفضة دون تغيير العتاد، لكنها مقيدة بالفجوة الكبيرة في السرعة بين Flash والذاكرة الفيزيائية، وليست “دواءً سحرياً”. هل تفعلها وكم حجمها، يجب الحكم بناءً على حجم الذاكرة الفيزيائية لطرازك، ومواصفات الـ Flash، ومساحة التخزين المتبقية، وعادات الاستخدام.

\end{tcbinfo}

\begin{tcbpractice}{لنبدأ عملياً: تحقق إن كان هاتفك لديه توسيع ذاكرة}
\begin{enumerate}
  \item افتح [الإعدادات]، ابحث عن "توسيع الذاكرة"، "دمج الذاكرة"، "RAM Plus" أو "Extended RAM".
  \item بعد إيجاده، تحقق مما إذا كان مفعلاً حالياً، وما هو إعداد سعة التوسيع.
  \item قارن مع محتوى هذا القسم: ما حجم الذاكرة الفيزيائية لطرازك؟ ما مواصفات الـ Flash؟ هل مساحة التخزين المتبقية كافية؟
  \item جرب إيقاف أو خفض سعة التوسيع، راقب تغيرات إبقاء الخلفية والحرارة لمدة أسبوع، لإيجاد الإعداد الأنسب لك.
\end{enumerate}

\end{tcbpractice}
